\section{Summary}\label{summary}

\texttt{rbfmvar} implements the Residual-Based Fully Modified Vector
Autoregression (RBFM-VAR) estimator of Chang (2000). The RBFM-VAR
procedure extends the Phillips (1995) FM-VAR approach to handle any
unknown mixture of I(0), I(1), and I(2) components without requiring
prior knowledge of the number or location of unit roots. This robustness
is achieved by applying a fully modified correction to VAR residuals
rather than to levels, yielding asymptotically chi-squared Wald
statistics for Granger non-causality tests regardless of integration
order. The package provides automatic lag selection via information
criteria (AIC, BIC, HQ), long-run variance estimation using Bartlett,
Parzen, or Quadratic Spectral kernels with Andrews (1991) automatic
bandwidth selection, impulse response functions (IRF) with bootstrap
confidence intervals, forecast error variance decomposition (FEVD), and
out-of-sample forecasting.

\section{Statement of Need}\label{statement-of-need}

Standard VAR inference breaks down when the system contains integrated
or cointegrated variables: OLS Wald statistics for Granger causality
tests have non-standard distributions, and pre-tests for integration
order are required before inference. The Toda-Yamamoto (1995) approach
augments the VAR lag length but requires knowledge of the maximum
integration order. The RBFM-VAR of Chang (2000) sidesteps both issues by
correcting for serial correlation and endogeneity in the residuals,
producing standard chi-squared inference without pre-testing. While
RBFM-VAR implementations exist in GAUSS and EViews, no R package has
provided this estimator. \texttt{rbfmvar} fills this gap, enabling
robust causal inference in multivariate time series with mixed
integration orders entirely within R.

\section{Usage}\label{usage}

\subsection{Fitting the RBFM-VAR
Model}\label{fitting-the-rbfm-var-model}

\begin{Shaded}
\begin{Highlighting}[]
\FunctionTok{library}\NormalTok{(rbfmvar)}

\CommentTok{\# Fit RBFM{-}VAR with automatic lag selection}
\NormalTok{fit }\OtherTok{\textless{}{-}} \FunctionTok{rbfmvar}\NormalTok{(}\AttributeTok{data =}\NormalTok{ macro\_ts, }\AttributeTok{lag\_max =} \DecValTok{4}\NormalTok{, }\AttributeTok{ic =} \StringTok{"BIC"}\NormalTok{,}
               \AttributeTok{kernel =} \StringTok{"Bartlett"}\NormalTok{)}
\FunctionTok{summary}\NormalTok{(fit)}
\FunctionTok{print}\NormalTok{(fit)}
\end{Highlighting}
\end{Shaded}

\subsection{Granger Non-Causality
Tests}\label{granger-non-causality-tests}

\begin{Shaded}
\begin{Highlighting}[]
\CommentTok{\# Wald test: does x Granger{-}cause y? (asymptotic chi{-}squared)}
\NormalTok{gc\_result }\OtherTok{\textless{}{-}} \FunctionTok{granger}\NormalTok{(fit, }\AttributeTok{cause =} \StringTok{"x"}\NormalTok{, }\AttributeTok{effect =} \StringTok{"y"}\NormalTok{)}
\FunctionTok{print}\NormalTok{(gc\_result)}
\end{Highlighting}
\end{Shaded}

\subsection{Impulse Response
Functions}\label{impulse-response-functions}

\begin{Shaded}
\begin{Highlighting}[]
\CommentTok{\# IRF with bootstrap confidence intervals}
\NormalTok{irf\_result }\OtherTok{\textless{}{-}} \FunctionTok{irf}\NormalTok{(fit, }\AttributeTok{impulse =} \StringTok{"x"}\NormalTok{, }\AttributeTok{response =} \StringTok{"y"}\NormalTok{,}
                  \AttributeTok{h =} \DecValTok{20}\NormalTok{, }\AttributeTok{nboot =} \DecValTok{1000}\NormalTok{, }\AttributeTok{ci =} \FloatTok{0.95}\NormalTok{)}
\FunctionTok{plot}\NormalTok{(irf\_result)}
\end{Highlighting}
\end{Shaded}

\subsection{Forecast Error Variance
Decomposition}\label{forecast-error-variance-decomposition}

\begin{Shaded}
\begin{Highlighting}[]
\NormalTok{fevd\_result }\OtherTok{\textless{}{-}} \FunctionTok{fevd}\NormalTok{(fit, }\AttributeTok{h =} \DecValTok{20}\NormalTok{)}
\FunctionTok{print}\NormalTok{(fevd\_result)}
\FunctionTok{plot}\NormalTok{(fevd\_result)}
\end{Highlighting}
\end{Shaded}

\subsection{Out-of-Sample Forecasting}\label{out-of-sample-forecasting}

\begin{Shaded}
\begin{Highlighting}[]
\NormalTok{fc }\OtherTok{\textless{}{-}} \FunctionTok{forecast}\NormalTok{(fit, }\AttributeTok{h =} \DecValTok{8}\NormalTok{)}
\FunctionTok{print}\NormalTok{(fc)}
\end{Highlighting}
\end{Shaded}

\section{Implementation}\label{implementation}

\texttt{rbfmvar} is implemented in R with a dependency on \texttt{MASS}
for matrix operations. Lag selection in \texttt{lag\_selection()}
evaluates AIC, BIC, and Hannan-Quinn criteria over a grid of candidate
orders. Long-run variance estimation in \texttt{lrv()} uses the
\texttt{kernels()} module which implements Bartlett, Parzen, and
Quadratic Spectral kernels with Andrews (1991) data-driven bandwidth.
The core \texttt{rbfmvar()} function estimates the VAR by OLS, extracts
residuals, applies the FM correction following Chang (2000), and
constructs the RBFM coefficient matrix. Granger causality Wald
statistics in \texttt{granger()} are asymptotically chi-squared under
the null regardless of integration order. IRFs are computed via the
moving-average representation and bootstrapped using a residual
bootstrap with 1000 replications by default. FEVD decomposition in
\texttt{fevd()} uses the Cholesky factorisation of the residual
covariance matrix.

\section*{References}\label{references}
\addcontentsline{toc}{section}{References}

\protect\phantomsection\label{refs}
\begin{CSLReferences}{1}{1}
\bibitem[\citeproctext]{ref-Andrews1991}
Andrews, Donald W. K. 1991. {``Heteroskedasticity and Autocorrelation
Consistent Covariance Matrix Estimation.''} \emph{Econometrica} 59 (3):
817--58. \url{https://doi.org/10.2307/2938229}.

\bibitem[\citeproctext]{ref-Chang2000}
Chang, Yoosoon. 2000. {``Bootstrap Unit Root Tests in Models with
{GARCH} Innovations.''} \emph{Econometric Theory} 16 (6): 787--818.
\url{https://doi.org/10.1017/S0266466600166071}.

\bibitem[\citeproctext]{ref-Phillips1995}
Phillips, Peter C. B. 1995. {``Fully Modified Least Squares and Vector
Autoregression.''} \emph{Econometrica} 63 (5): 1023--78.
\url{https://doi.org/10.2307/2171721}.

\bibitem[\citeproctext]{ref-TodaYamamoto1995}
Toda, Hiro Y., and Taku Yamamoto. 1995. {``Statistical Inference in
Vector Autoregressions with Possibly Integrated Processes.''}
\emph{Journal of Econometrics} 66 (1-2): 225--50.
\url{https://doi.org/10.1016/0304-4076(94)01616-8}.

\end{CSLReferences}
